\appendices

\section{Computer Backend (Design, Not Implemented)}
\label{app:computer}

The repository contains a desktop executor that wraps PyAutoGUI and psutil, a text-file helper and a file-existence predicate. None is connected to the control loop: no node, action type or API route calls them, and the file helper does not confine paths. The rest of this appendix is a design.

\paragraph{Observation and grounding} The backend would observe the foreground window through its accessibility tree (Windows UI Automation, macOS Accessibility, AT-SPI) into a manifest shaped like the DOM manifest, falling back to a vision model on a window screenshot. File and process actions are grounded in typed state (paths, hashes, process identity).

\paragraph{Postconditions and permission tiers} Each mutating action carries exact postconditions (Table~\ref{tab:post}) captured with before/after values. Because many local actions cannot be undone, a postcondition checked afterwards is no protection on its own, so an \texttt{action\_gate} node between planning and execution would compute a tier from the concrete action: T0 (read-only) runs freely; T1 (reversible, with a recorded inverse) runs automatically only if the user enabled it; T2 (destructive) and T3 (network-capable programs, credentials) always need approval bound to a hash of the exact action. Paths outside allow-listed roots, inline interpreters and privilege elevation would be refused. A design in which untrusted content cannot change the control flow~\cite{debenedetti2025camel} is a stronger alternative.

\begin{table}[t]
\centering
\caption{Completion postconditions by environment and their status in the code. Impl.: implemented and used by the agent; Unit: implemented and unit-tested, not called by the agent; Prop.: proposed design. $h_0$: hash before the action; $(\mathit{pid},t_c)$: process ID and creation time.}
\label{tab:post}
\footnotesize
\setlength{\tabcolsep}{2.2pt}
\begin{tabular}{@{}p{1.25cm}p{5.05cm}l@{}}
\toprule
Env. & Example postcondition & Status \\
\midrule
Browser & no error or sign-in page; results page for the quoted query; requested text present; answer mentions the goal; no pending step & Impl. \\
Browser & bot-check page detected $\Rightarrow$ \textsc{Blocked} & Impl. \\
File & path exists and size $\geq n$ & Unit \\
File & created / overwritten: $\mathrm{sha256}=\mathrm{sha256}(\text{content})$; moved: source absent and $\mathrm{sha256}(\text{dst})=h_0(\text{src})$; trashed: path absent and trash entry present & Prop. \\
Process & launched: $(\mathit{pid},t_c)$ alive; terminated: $(\mathit{pid},t_c)$ gone within a timeout & Prop. \\
Window & foreground title matches; window with title exists & Prop. \\
Desktop UI & accessibility value of a control equals the input; control state contains \textit{checked} & Prop. \\
Read-only & workspace manifest (paths, sizes, mtimes) unchanged & Prop. \\
\bottomrule
\end{tabular}
\end{table}


\section{End-to-End Task Set}
\label{app:tasks}

The 40 tasks (\texttt{benchmarks/e2e/tasks.json}) are distributed as follows (development/test): extraction 2/8, form submission 2/7, multi-step 1/5, navigation 2/2, search 1/2, text entry 1/2, unsatisfiable 1/4; tiers 1--4 contain 3/6, 4/11, 3/12 and 0/1 tasks. The test split uses all five sites; the notes site appears only there. The harness (\texttt{run\_e2e.py}) takes the split, the number of trials and a list of the configurations of Table~\ref{tab:e2e}; the analysis script (\texttt{analyze\_e2e.py}) writes result tables with Wilson intervals, McNemar tests and bootstrap intervals.
